\section{Additional Experimental Results}
\subsection{Overall Performance on RTX 5090}
To validate the generality of our approach across GPU architectures, we conduct comprehensive experiments on NVIDIA GeForce RTX 5090, which features the Blackwell architecture with 170 SMs, 21,760 CUDA cores, and 32GB GDDR6X memory. Table~\ref{tab:5090} presents the results, which closely match the RTX 4090 performance reported in the main text (success rate=100\%, average speedup$>$2), demonstrating the strong generalization capability and architectural robustness of our framework.

\begin{table}[h]
\centering
\caption{Performance summary on RTX 5090 across KernelBench difficulty levels.}
\label{tab:5090}
\begin{tabular}{lccc}
\toprule
\textbf{Level} & \textbf{Success Rate (\%)} & \textbf{$\mathbf{Fast}_1$ Rate (\%)} & \textbf{Avg. Speedup ($\times$)} \\
\midrule
Level 1 (Basic)      & 100   & 69 & 2.40 \\
Level 2 (Medium)     & 100   & 96 & 2.67 \\
Level 3 (Hard)       & 100   & 90 & 1.41 \\
\midrule
\textbf{Overall}     & \textbf{100} & \textbf{84} & \textbf{2.31} \\
\bottomrule
\end{tabular}
\end{table}

\section{Prompt Details}

This section provides the complete prompts used in our multi-agent framework. These prompts encode domain-specific optimization knowledge and guide agents through different stages of kernel generation.

\subsection{Seed Prompt}

The seed prompt is used during the multi-seed search stage to generate initial Triton kernel implementations. It receives the target PyTorch operator (\texttt{\$kernel\_src}) and few-shot examples (\texttt{\$few\_base} and \texttt{\$few\_new}) as input. The prompt emphasizes strict syntactic constraints and common pitfalls to ensure high initial correctness. It outputs initial Triton kernel implementations with diverse algorithmic structures.

\begin{mdframed}[linewidth=1pt, linecolor=black, backgroundcolor=gray!5]
\begin{lstlisting}
Write high-performance Triton kernels to replace PyTorch operators.
Generate the FASTEST kernel while maintaining correctness.

## CRITICAL — These cause 60%+ of failures:
1. EVERY kernel function MUST have `@triton.jit` decorator — MANDATORY
2. Grid size MUST be > 0: use `triton.cdiv(N, BLOCK)` or `max(1, N // BLOCK)`
3. BLOCK sizes MUST be power-of-2 constexpr: 16, 32, 64, 128, 256
4. `tl.program_id(axis)` only supports axis = 0, 1, 2 (max 3D grid)

## Triton Syntax Rules:
- For matmul/conv/linear ops, prefer `tl.dot(a, b, allow_tf32=True)` over element-wise multiply-add
- No `continue`, `break`, `return` inside loops — use masking instead
- No tensor indexing with loop vars: `x[:, i]` or `x[i, :]` is INVALID
- No tuple unpacking inside kernel: `a, b = tl.load(...)` is INVALID
- No nested functions inside @triton.jit
- No Python control flow on tl.tensor or BLOCK_* values
- No dynamic `tl.reshape()` or view operations

## Missing Triton Functions (implement manually):
- tl.tanh → `(tl.exp(2*x) - 1) / (tl.exp(2*x) + 1)`
- tl.sigmoid → `1 / (1 + tl.exp(-x))`
- tl.gelu, tl.silu, tl.softmax, tl.mish → implement from definition

## Load/Store Rules:
- Pointer + scalar offset → scalar value
- Pointer + block offset (via tl.arange) → block of values
- mask shape MUST match data shape exactly

## Output Format (STRICT):
1. Imports: `import torch, torch.nn as nn, triton, triton.language as tl` (and math if needed)
2. `@triton.jit` kernel(s) — MUST have this decorator
3. Wrapper function with grid calculation
4. `class ModelNew(nn.Module)` — REQUIRED

Do NOT include: testing code, `if __name__ == "__main__"`, get_inputs, get_init_inputs

Example PyTorch:
'''
$few_base
'''

Example Triton:
'''
$few_new
'''

Target:
```python
$kernel_src
"""
\end{lstlisting}
\end{mdframed}

\subsection{Stage System Prompts for Hardware-Specific Tuning}

During hardware-specific tuning, each optimization stage uses a specialized system prompt that defines the stage's focus, relevant metrics, and optimization rules. These prompts constrain agent decisions to stage-specific objectives, preventing cross-stage interference.

\subsubsection{Grid and Parallel Mapping}

This stage focuses on determining the optimal mapping between operator dimensions and GPU grid dimensions. The prompt guides the agent to prioritize batch/head/expert parallelism before reducing block sizes, and ensures grid configurations remain within hardware limits (maximum 3 dimensions).

\begin{mdframed}[linewidth=1pt, linecolor=black, backgroundcolor=gray!5]
\begin{lstlisting}
Focus: Grid layout & parallelism.

Metrics:
- sm__throughput.avg.pct_of_peak_sustained_elapsed (>60%)
- launch__grid_size

Rules:
- 1D: (cdiv(N, BLOCK))
- 2D: (cdiv(M, BLOCK_M), cdiv(N, BLOCK_N))
- 3D: (batch, cdiv(M, BLOCK_M), cdiv(N, BLOCK_N))
- >3D: flatten ONLY independent dims
- Prefer batch / head / expert / group parallelism before shrinking BLOCK
- For grouped operations: ensure group dimension is in grid (e.g., program_id(2) for groups)
- Change grid only if SM utilization is clearly low

Safety:
- Max 3 grid dims, static rank
- grid=(G0,G1,G2) must match tl.program_id(0/1/2)
- For grouped ops: verify group indexing is correct
- If unsure about correctness, do NOT change grid

Autotune:
- Autotune either BLOCK_* OR (num_warps, num_stages)
- If autotuning BLOCK_*, use grid=lambda META: (...)
- Never redefine BLOCK_* in both kernel and launch
- Max 2-3 configs to reduce compilation time
\end{lstlisting}
\end{mdframed}

\subsubsection{Tensor Tiling}

This stage optimizes the granularity of computation within each thread block by selecting appropriate \texttt{BLOCK\_M}, \texttt{BLOCK\_N}, and \texttt{BLOCK\_K} sizes. The prompt enforces power-of-2 constraints and guides autotuning across a small set of configurations to balance data reuse with register pressure.

\begin{mdframed}[linewidth=1pt, linecolor=black, backgroundcolor=gray!5]
\begin{lstlisting}
Focus: BLOCK_M/N/K selection.

Metrics:
- sm__warps_active.avg.pct_of_peak_sustained_active (>50%)

Rules:
- BLOCK_* must be powers of 2
- Tensor Core: BLOCK_M/N multiple of 16, BLOCK_K multiple of 8 (preference)
- FP32: M/N ∈ {32,64,128,256}, K ∈ {16,32,64}
- Avoid oversized tiles (mask waste)
- Keep baseline tile if unsure

Autotune:
- Max 2-3 configs to reduce compilation time
- Autotune ONLY on @triton.jit kernel
""",
\end{lstlisting}
\end{mdframed}

\subsubsection{Memory and Tuning}

The final stage refines memory access patterns and pipeline execution. The prompt directs the agent to tune \texttt{num\_warps} and \texttt{num\_stages} based on occupancy and memory stall metrics, while keeping grid configuration and block sizes fixed from previous stages.

\begin{mdframed}[linewidth=1pt, linecolor=black, backgroundcolor=gray!5]
\begin{lstlisting}
Focus: Memory optimization and final parameter tuning.

Metrics:
- dram__throughput.avg.pct_of_peak_sustained_elapsed
- lts__t_sector_hit_rate.pct
- smsp__warp_issue_stalled_memory_dependency_per_warp_active.pct (<20%)
- sm__warps_active.avg.pct_of_peak_sustained_active

Parameters to tune:
- num_stages ∈ {2, 3, 4}
- num_warps ∈ {4, 8} (based on occupancy)

Rules:
- Increase num_stages only if memory stalls > 20%
- Change num_warps only if occupancy suggests it
- Larger BLOCK_K improves reuse but increases register pressure
- Do NOT modify grid or BLOCK sizes (fixed in earlier stages)
- Do not rewrite access patterns without metric evidence

Autotune:
- Max 3-4 configs combining num_stages and num_warps
- Always include original config as baseline
- Revert if gain < 2% or unstable
"""
\end{lstlisting}
\end{mdframed}

\subsection{Profile Agent Prompts}

The profile agent analyzes kernel performance and identifies optimization opportunities. At different stages, it receives different system prompts to align with stage-specific objectives.

\subsubsection{Algorithmic Refinement Stage}

During algorithmic refinement, the profile agent receives the PyTorch reference code, current Triton kernel implementation, and performance metrics as input. The prompt guides the agent to analyze high-level algorithmic optimizations (operator fusion, algorithm replacement, etc.) through structured analysis: code inspection, performance diagnosis, and root cause identification. It outputs a JSON response specifying the identified bottleneck, the proposed optimization method, and the modification plan for structural transformation.


\begin{mdframed}[linewidth=1pt, linecolor=black, backgroundcolor=gray!5]
\begin{lstlisting}

You are a GPU kernel optimization architect. Analyze the kernel and identify
**ONE high-level algorithmic optimization**.

# PyTorch Reference
```python
$python_code
```

# Current Triton Kernel
```python
$triton_code
```

$performance_section

---

## Analysis Steps

1. **Code Analysis**: Count kernels, identify operations, check for inefficiencies
2. **Performance Diagnosis**: Use metrics/latency to identify bottleneck type
3. **Root Cause**: Combine code + performance to find the core issue

## Optimization Categories (pick ONE if worth optimizing):

### 1. Operator Fusion
Fuse consecutive ops into fewer kernels to reduce memory traffic and launch overhead.

### 2. Algorithm Replacement
Replace naive algorithm with optimized variant.
- For Attention: Flash Attention, online softmax
- For Convolution: Winograd, im2col
- For RNN/GRU/LSTM: Persistent kernel with HYBRID computation

### 3. Kernel Launch Reduction
Combine multiple small kernels to reduce overhead.

### 4. Memory Layout Optimization
Use in-place operations, buffer reuse, or better layouts.

## Should We Optimize?

Before proposing optimization, determine if it's worthwhile:
- **Not worth optimizing** if:
  - Code is already near-optimal (expected speedup < 10%)
  - Bottleneck cannot be addressed (hardware limited, already optimal algorithm)
  - Optimization would add significant complexity with minimal gain

- **Worth optimizing** if:
  - Clear algorithmic inefficiency exists (multiple kernels, suboptimal algorithm)
  - Expected speedup >= 20%
  - Concrete optimization path available

## Output (JSON)

```json
{
  "worth_optimizing": "yes/no",
  "bottleneck": "<Root cause in 1-2 sentences>",
  "optimization method": "<Specific optimization in 1-2 sentences>",
  "modification plan": "<Implementation steps in 2-3 sentences>",
}
```

Return JSON only.
\end{lstlisting}
\end{mdframed}

\subsubsection{Hardware-Specific Tuning Stages}

During hardware-specific tuning, the profile agent receives NCU profiling metrics, current kernel code, and stage-specific system prompts as input. The prompt constrains analysis to the current stage's optimization scope (e.g., grid configuration, block sizes, or memory patterns) and directs the agent to propose exactly one targeted modification based on measured bottlenecks. It outputs a JSON response specifying the bottleneck, optimization method, and modification plan.

\begin{mdframed}[linewidth=1pt, linecolor=black, backgroundcolor=gray!5]
\begin{lstlisting}
You are a senior Triton kernel optimization engineer. Read the PyTorch reference code, the current Triton candidate, and the Nsight Compute metrics. Then identify **exactly one** highest-impact speed bottleneck, propose **exactly one** optimization method and propose a modification plan. Be surgical and metrics-driven.

# PyTorch Reference
$python_code

# Current Triton Kernel
```python
$TRITON_CODE
```

$STAGE_CONTEXT

# Nsight Compute Metrics (3 Core Metrics)
$NCU_METRICS

Rules:
- Return **one and only one** optimization method — the largest expected speedup.
- Focus on Triton-specific optimizations:
  * **BLOCK_M/N/K tuning**: Adjust tile sizes to optimize data reuse and cache efficiency
  * **num_warps**: Control occupancy (2/4/8 warps per block)
  * **num_stages**: Enable software pipelining (2-4 stages for memory-bound kernels)
  * **Memory access patterns**: Optimize coalescing, use tl.trans() for layout changes
  * **Grid configuration**: Adjust program_id mapping and workload distribution
- Prefer changes that directly address measured bottlenecks from NCU metrics:
  * High DRAM throughput → Increase BLOCK size for data reuse
  * Low cache hit rate → Adjust BLOCK size for better locality
  * Low occupancy → Tune num_warps, reduce register pressure
- Keep fields brief; avoid lists of alternatives, disclaimers, or generic advice.

Output format (JSON):
```json
{
  "bottleneck": "<max 30 words>",
  "optimization method": "<max 35 words>",
  "modification plan": "<max 35 words>"
}```
Read everything and follow the Rules exactly. Return the JSON in the specified format.
"""
\end{lstlisting}
\end{mdframed}

\subsection{Code Agent Prompts}

The code agent generates or modifies kernel implementations based on optimization suggestions or error diagnostics. Different prompts guide repair and generation tasks.

\subsubsection{Repair Mode}

When kernel execution fails, the code agent receives error logs and debug analysis as input. The prompt emphasizes strict adherence to Triton syntax rules and output format requirements. It outputs corrected Triton kernel code.

\begin{mdframed}[linewidth=1pt, linecolor=black, backgroundcolor=gray!5]
\begin{lstlisting}
Fix the Triton kernel errors. Generate correct code.

## ERROR LOG
```
$ERROR_LOG
```
## Broken Code
```python
$OLD_CODE
```

## ANALYSIS
```
$ANALYSIS_INFO
```

## OUTPUT FORMAT (STRICT):
1. Imports: torch, torch.nn, triton, triton.language as tl (and math if needed)
2. @triton.jit decorated kernel function(s)
3. Wrapper function(s) for grid calculation and kernel launch
4. class ModelNew(nn.Module) — REQUIRED

Do NOT include: testing code, if __name__, get_inputs, get_init_inputs

```python
# <corrected code>
```
\end{lstlisting}
\end{mdframed}

\subsubsection{Generation Mode (Algorithmic Refinement)}

During algorithmic refinement, the code agent receives the current kernel, PyTorch reference, and optimization analysis as input. The prompt directs the agent to implement structural transformations suggested by the profile agent and apply specific algorithmic changes while preserving correctness. It outputs the optimized Triton kernel code.

\begin{mdframed}[linewidth=1pt, linecolor=black, backgroundcolor=gray!5]
\begin{lstlisting}
You are optimizing a Triton kernel based on algorithmic analysis.

# PyTorch Reference (Target Behavior)

```python
$pytorch_reference
```

# ANALYSIS
```
$ANALYSIS_INFO
```

---

# Current Kernel (needs optimization)

```python
$current_kernel
```

---

# Your Task

Implement the optimization strategy above. Focus on the specific bottleneck identified.

## Key Requirements

1. **Preserve correctness**: Maintain the same input/output behavior
2. **Apply the optimization**: Follow the implementation plan exactly
3. **Use valid Triton syntax**:
   - Every kernel MUST have `@triton.jit` decorator
   - Grid size MUST be > 0: use `triton.cdiv(N, BLOCK)` or `max(1, N // BLOCK)`
   - BLOCK sizes MUST be power-of-2: 16, 32, 64, 128, 256
   - No `continue`, `break`, `return` inside kernels (use masking)
   - Prefer `tl.dot(a, b, allow_tf32=True)` for matmul operations
4. **Output format**:
   - Imports: `import torch, torch.nn as nn, triton, triton.language as tl`
   - `@triton.jit` kernel(s)
   - Wrapper function(s)
   - `class ModelNew(nn.Module)` — REQUIRED
   - NO testing code, NO `if __name__ == "__main__"`

---

Generate the optimized kernel now. Output ONLY the complete Python code.
"""
\end{lstlisting}
\end{mdframed}

\subsubsection{Generation Mode (Hardware-Specific Tuning)}

During hardware-specific tuning stages, the code agent receives the current kernel, target GPU information, NCU profiling metrics, stage-specific system prompts, and optimization analysis as input. The prompt directs the agent to apply focused hardware-level optimizations based on the profile agent's suggestions. It outputs the optimized Triton kernel code.

\begin{mdframed}[linewidth=1pt, linecolor=black, backgroundcolor=gray!5]
\begin{lstlisting}
You are a Triton kernel optimization specialist. Generate the FASTEST possible kernel.

# Target GPU: $gpu_name

[OPTIMIZATION STAGE]
$STAGE_CONTEXT

[CURRENT CODE]
```python
$arch_src
```

[ANALYSIS]
```
$ANALYSIS_INFO
```

**Task**: Analyze the NCU metrics and current code, then generate optimized code that maximizes performance.

## CRITICAL — Code MUST compile and run:
1. EVERY kernel function MUST have `@triton.jit` decorator
2. Grid size MUST be > 0: use `triton.cdiv(N, BLOCK)` or `max(1, N // BLOCK)`
3. BLOCK sizes MUST be power-of-2: 16, 32, 64, 128, 256
4. `tl.program_id(axis)` only supports axis = 0, 1, 2
5. No `continue`, `break`, `return` inside loops — use masking
6. No tensor indexing with loop vars: `x[:, i]` is INVALID
7. mask shape MUST match data shape in tl.load/tl.store

## Missing Triton Functions (implement manually):
- tl.tanh, tl.sigmoid, tl.gelu, tl.silu, tl.softmax, tl.mish

## OUTPUT FORMAT (STRICT):
1. Imports: torch, torch.nn, triton, triton.language as tl
2. @triton.jit decorated kernel function(s)
3. Wrapper function(s) for grid calculation and kernel launch
4. class ModelNew(nn.Module) that calls your kernels

Do NOT include: testing code, if __name__, get_inputs, get_init_inputs

```python
# <optimized Triton code>
```
\end{lstlisting}
\end{mdframed}

\subsection{Debug Agent Prompt}

The debug agent receives error logs, PyTorch reference implementations, and broken kernel code as input. The prompt guides the agent to diagnose execution failures and identify root causes with specific, actionable diagnostic information. It outputs a structured JSON response containing the critical issue, its impact, and required modifications, which the code agent then uses to generate fixes.

\begin{mdframed}[linewidth=1pt, linecolor=black, backgroundcolor=gray!5]
\begin{lstlisting}
You are a Triton kernel debugging expert. Analyze the error and identify the root cause.

## ERROR LOG
```
$ERROR_LOG
```

## Expected Behavior (PyTorch Reference)
```python
$PYTORCH_CODE
```

## Current Implementation (Broken Triton Kernel)
```python
$CUDA_CODE
```

---

## Your Task

Identify the **single most critical issue** that causes the error above.

### Analysis Guidelines

1. **Focus on root cause**, not symptoms
   - Bad: "Output is wrong"
   - Good: "BLOCK_K loop missing, only processes first 32 elements of K dimension"

2. **Be specific about WHAT and WHERE**
   - Bad: "Memory access issue"
   - Good: "Line 45: tl.atomic_add(c_block_ptr, acc) - atomic_add requires scalar pointer, not block_ptr"

3. **Prioritize by impact**
   - Correctness bugs > Performance issues > Style problems
   - Algorithm errors > Implementation details

### Output Format

**CRITICAL: You MUST output ONLY valid JSON. No other text allowed.**

```json
{
  "critical_issue": "<Concise description of THE root cause, max 30 words>",
  "why_it_matters": "<Why this causes the observed error, max 35 words>",
  "modification plan": "<What needs to change (not how), max 30 words>"
}
```
**Remember**: Output ONLY the JSON block. No explanations, no commentary, no additional text.
\end{lstlisting}
\end{mdframed}

\section{Nsight Compute Profiling Metrics}

NVIDIA Nsight Compute (NCU) provides detailed hardware-level performance metrics that guide optimization decisions at each stage. This section documents the metrics we collect and how they inform stage-specific optimization strategies.

\subsection{Complete Metric Set}

We collect a comprehensive set of metrics covering compute utilization, memory hierarchy performance, and execution bottlenecks. The following list shows all metrics gathered during profiling:

\begin{mdframed}[linewidth=1pt, linecolor=black, backgroundcolor=gray!5]
\begin{lstlisting}
    "sm__cycles_active.avg",
    "sm__warps_active.avg.pct_of_peak_sustained_active",
    "launch__occupancy_limit_blocks",
    "launch__occupancy_limit_registers",
    "launch__occupancy_limit_shared_mem",
    "launch__registers_per_thread",
    "sm__inst_executed.sum",
    "sm__inst_executed_pipe_fp32.avg.pct_of_peak_sustained_active",
    "sm__inst_executed_pipe_tensor.avg.pct_of_peak_sustained_active",
    "dram__bytes_read.sum",
    "dram__bytes_write.sum",
    "dram__throughput.avg.pct_of_peak_sustained_elapsed",
    "dram__bytes.sum.per_second",
    "gpu__dram_throughput.avg.pct_of_peak_sustained_elapsed",
    "l1tex__t_sector_hit_rate.pct",
    "l1tex__throughput.avg.pct_of_peak_sustained_active",
    "lts__t_sector_hit_rate.pct",
    "lts__throughput.avg.pct_of_peak_sustained_active",
    "smsp__warp_issue_stalled_memory_dependency_per_warp_active.pct",
    "smsp__warp_issue_stalled_short_scoreboard_per_warp_active.pct",
    "smsp__warp_issue_stalled_long_scoreboard_per_warp_active.pct",
    "smsp__warp_issue_stalled_barrier_per_warp_active.pct",
    "smsp__warp_issue_stalled_branch_resolving_per_warp_active.pct",
    "smsp__sass_average_branch_targets_threads_uniform.pct"
\end{lstlisting}
\end{mdframed}

These metrics span multiple aspects of GPU execution:
\begin{itemize}
\item \textbf{Compute utilization:} SM throughput, active warps, instruction execution rates
\item \textbf{Memory hierarchy:} DRAM bandwidth, L1/L2 cache hit rates, memory throughput
\item \textbf{Execution stalls:} Memory dependency stalls, scoreboard stalls, barrier stalls
\item \textbf{Resource limits:} Occupancy limits due to blocks/registers/shared memory
\item \textbf{Control flow:} Branch uniformity and divergence patterns
\end{itemize}

\subsection{Core Metrics}

A subset of core metrics provides consistent baseline performance indicators across all optimization stages. These metrics are always collected to enable cross-stage performance comparison:

\begin{mdframed}[linewidth=1pt, linecolor=black, backgroundcolor=gray!5]
\begin{lstlisting}
"sm__throughput.avg.pct_of_peak_sustained_elapsed",  # SM compute utilization
"dram__throughput.avg.pct_of_peak_sustained_elapsed",  # Memory bandwidth
"lts__t_sector_hit_rate.pct",  # L2 cache hit rate
"sm__warps_active.avg.pct_of_peak_sustained_active",  # Warp occupancy
"smsp__sass_average_data_bytes_per_sector_mem_global_op_ld.pct",  # Memory coalescing
\end{lstlisting}
\end{mdframed}

\subsection{Stage-Specific Metrics}

Each hardware-specific tuning stage focuses on a targeted subset of metrics that directly inform the stage's optimization objectives.

\subsubsection{Grid and Parallel Mapping Metrics}

These metrics evaluate how effectively computation is distributed across GPU cores and whether additional parallelism can be exploited:

\begin{mdframed}[linewidth=1pt, linecolor=black, backgroundcolor=gray!5]
\begin{lstlisting}
"launch__grid_size",
"launch__block_size",
"launch__waves_per_multiprocessor"
\end{lstlisting}
\end{mdframed}

Low \texttt{launch\_\_grid\_size} or \texttt{waves\_per\_multiprocessor} suggests insufficient parallelism, indicating opportunities to increase grid dimensions or adjust block-level parallelism.

\subsubsection{Block Tiling Metrics}

These metrics diagnose whether block size choices create resource bottlenecks that limit occupancy:

\begin{mdframed}[linewidth=1pt, linecolor=black, backgroundcolor=gray!5]
\begin{lstlisting}
"launch__occupancy_limit_blocks",
"launch__occupancy_limit_registers",
"launch__occupancy_limit_shared_mem",
"launch__registers_per_thread",
"launch__shared_mem_per_block_static"
\end{lstlisting}
\end{mdframed}

The \texttt{occupancy\_limit} metrics reveal whether register usage, shared memory usage, or thread block count constrains the number of active warps. This guides adjustments to \texttt{BLOCK\_M/N/K} sizes to balance data reuse with resource consumption.

\subsubsection{Memory and Tuning Metrics}

These metrics identify memory access inefficiencies and pipeline stalls that can be addressed through software pipelining and access pattern optimization:

\begin{mdframed}[linewidth=1pt, linecolor=black, backgroundcolor=gray!5]
\begin{lstlisting}
"smsp__warp_issue_stalled_memory_dependency_per_warp_active.pct",
"l1tex__t_sector_hit_rate.pct",
"smsp__sass_average_data_bytes_per_sector_mem_global_op_st.pct",
"smsp__warp_issue_stalled_short_scoreboard_per_warp_active.pct",
"smsp__warp_issue_stalled_barrier_per_warp_active.pct",
"sm__inst_executed_pipe_tensor.avg.pct_of_peak_sustained_active"
\end{lstlisting}
\end{mdframed}

High \texttt{warp\_issue\_stalled\_memory\_dependency} percentages indicate memory latency bottlenecks that can benefit from increasing \texttt{num\_stages} for software pipelining. Low cache hit rates and poor memory coalescing (\texttt{bytes\_per\_sector}) suggest access pattern issues. The \texttt{inst\_executed\_pipe\_tensor} metric indicates Tensor Core utilization for matrix operations.

By providing stage-specific metric subsets to the profile agent, we constrain its analysis to the current optimization objective, enabling focused and effective optimization recommendations.

\section{Benchmark Details}

This section provides detailed specifications of the benchmarks used in our evaluation, including representative task examples with their PyTorch reference implementations.

\subsection{KernelBench}

KernelBench consists of 250 tasks spanning three difficulty levels. Each task provides a standardized PyTorch reference implementation following a consistent interface: a \texttt{Model} class with forward method, \texttt{get\_inputs()} for runtime inputs, and \texttt{get\_init\_inputs()} for model initialization parameters. Below we present one representative example from each difficulty level.

\subsubsection{Level 1: Square Matrix Multiplication}
\begin{mdframed}[linewidth=1pt, linecolor=black, backgroundcolor=gray!5]
\begin{lstlisting}
import torch
import torch.nn as nn

class Model(nn.Module):
    """
    Simple model that performs a single square matrix multiplication (C = A * B)
    """
    def __init__(self):
        super(Model, self).__init__()
    
    def forward(self, A: torch.Tensor, B: torch.Tensor) -> torch.Tensor:
        """
        Performs the matrix multiplication.

        Args:
            A (torch.Tensor): Input matrix A of shape (N, N).
            B (torch.Tensor): Input matrix B of shape (N, N).

        Returns:
            torch.Tensor: Output matrix C of shape (N, N).
        """
        return torch.matmul(A, B)

N = 2048 * 2

def get_inputs():
    A = torch.rand(N, N)
    B = torch.rand(N, N)
    return [A, B]

def get_init_inputs():
    return []  # No special initialization inputs needed
\end{lstlisting}
\end{mdframed}

\subsubsection{Level 2: Conv2D + ReLU + BiasAdd}
\begin{mdframed}[linewidth=1pt, linecolor=black, backgroundcolor=gray!5]
\begin{lstlisting}
import torch
import torch.nn as nn

class Model(nn.Module):
    """
    Simple model that performs a convolution, applies ReLU, and adds a bias term.
    """
    def __init__(self, in_channels, out_channels, kernel_size, bias_shape):
        super(Model, self).__init__()
        self.conv = nn.Conv2d(in_channels, out_channels, kernel_size)
        self.bias = nn.Parameter(torch.randn(bias_shape)) 

    def forward(self, x):
        x = self.conv(x)
        x = torch.relu(x)
        x = x + self.bias
        return x

batch_size = 128
in_channels  = 64  
out_channels = 128  
height = width = 128
kernel_size = 3
bias_shape = (out_channels, 1, 1)

def get_inputs():
    return [torch.rand(batch_size, in_channels, height, width)]

def get_init_inputs():
    return [in_channels, out_channels, kernel_size, bias_shape]
\end{lstlisting}
\end{mdframed}

\subsubsection{Level 3: Multi-Layer Perceptron (MLP)}
\begin{mdframed}[linewidth=1pt, linecolor=black, backgroundcolor=gray!5]
\begin{lstlisting}
import torch
import torch.nn as nn
import torch.nn.functional as F

class Model(nn.Module):
    def __init__(self, input_size, layer_sizes, output_size):
        """
        :param input_size: The number of input features
        :param layer_sizes: A list of ints containing the sizes of each hidden layer
        :param output_size: The number of output features
        """
        super(Model, self).__init__()
        
        layers = []
        current_input_size = input_size
        
        for layer_size in layer_sizes:
            layers.append(nn.Linear(current_input_size, layer_size))
            layers.append(nn.ReLU())
            current_input_size = layer_size
        
        layers.append(nn.Linear(current_input_size, output_size))
        
        self.network = nn.Sequential(*layers)
    
    def forward(self, x):
        """
        :param x: The input tensor, shape (batch_size, input_size)
        :return: The output tensor, shape (batch_size, output_size)
        """
        return self.network(x)

# Test code
batch_size = 128
input_size = 16384
layer_sizes = [16384, 16384]
output_size = 8192

def get_inputs():
    return [torch.rand(batch_size, input_size)]

def get_init_inputs():
    return [input_size, layer_sizes, output_size]
\end{lstlisting}
\end{mdframed}

\subsection{TritonBench}

TritonBench evaluates performance on production-grade kernels commonly used in LLM inference workloads. These operators are sourced from real-world Triton implementations in public repositories (e.g., turbo-llm/turbo-alignment). Since the original benchmarks provide natural language task descriptions, we converted them to standardized PyTorch reference implementations to align with our evaluation framework. Below we show an example RoPE (Rotary Position Embedding) operator, presenting both the original Triton implementation from GitHub and our converted PyTorch reference.

\subsubsection{Example: RoPE (Rotary Position Embedding)}

\textbf{Original Triton Implementation (from turbo-llm/turbo-alignment):}
\begin{mdframed}[linewidth=1pt, linecolor=black, backgroundcolor=gray!5]
\begin{lstlisting}

import torch
import triton
import triton.language as tl

@triton.jit
def _triton_rope(
    q_ptr,
    q_row_stride,
    k_ptr,
    k_row_stride,
    cos,
    cos_row_stride,
    sin,
    sin_row_stride,
    sl,
    bs: tl.constexpr,
    n_qh: tl.constexpr,
    n_kh: tl.constexpr,
    hd: tl.constexpr,
    pad_n_qh: tl.constexpr,
    pad_n_kh: tl.constexpr,
    pad_hd: tl.constexpr,
    BLOCK_SIZE: tl.constexpr,
    BACKWARD_PASS: tl.constexpr = False,
):
    pid = tl.program_id(0)

    q_ptr = q_ptr + pid * q_row_stride
    k_ptr = k_ptr + pid * k_row_stride

    cos_row_idx = pid % (sl)
    cos = cos + cos_row_idx * cos_row_stride
    sin = sin + cos_row_idx * sin_row_stride
    cos_offsets = tl.arange(0, pad_hd // 2)
    cos_mask = cos_offsets < hd // 2
    cos_row = tl.load(cos + cos_offsets, mask=cos_mask, other=0)
    sin_row = tl.load(sin + cos_offsets, mask=cos_mask, other=0)

    first_half_q_offsets = tl.arange(0, pad_n_qh)[:, None] * hd + tl.arange(0, pad_hd // 2)[None, :]
    first_half_k_offsets = tl.arange(0, pad_n_kh)[:, None] * hd + tl.arange(0, pad_hd // 2)[None, :]
    first_q_mask = (tl.arange(0, pad_n_qh)[:, None] < n_qh) & (tl.arange(0, pad_hd // 2)[None, :] < hd // 2)
    first_k_mask = (tl.arange(0, pad_n_kh)[:, None] < n_kh) & (tl.arange(0, pad_hd // 2)[None, :] < hd // 2)
    q_tile_1 = tl.load(q_ptr + first_half_q_offsets, mask=first_q_mask, other=0).to(sin_row.dtype)
    k_tile_1 = tl.load(k_ptr + first_half_k_offsets, mask=first_k_mask, other=0).to(sin_row.dtype)

    second_half_q_offsets = first_half_q_offsets + (hd // 2)
    second_half_k_offsets = first_half_k_offsets + (hd // 2)
    second_q_mask = first_q_mask
    second_k_mask = first_k_mask
    q_tile_2 = tl.load(q_ptr + second_half_q_offsets, mask=second_q_mask, other=0).to(sin_row.dtype)
    k_tile_2 = tl.load(k_ptr + second_half_k_offsets, mask=second_k_mask, other=0).to(sin_row.dtype)

    if not BACKWARD_PASS:
        new_q_tile_1 = q_tile_1 * cos_row - q_tile_2 * sin_row
        tl.store(q_ptr + first_half_q_offsets, new_q_tile_1, mask=first_q_mask)
        new_q_tile_2 = q_tile_2 * cos_row + q_tile_1 * sin_row
        tl.store(q_ptr + second_half_q_offsets, new_q_tile_2, mask=second_q_mask)

        new_k_tile_1 = k_tile_1 * cos_row - k_tile_2 * sin_row
        tl.store(k_ptr + first_half_k_offsets, new_k_tile_1, mask=first_k_mask)
        new_k_tile_2 = k_tile_2 * cos_row + k_tile_1 * sin_row
        tl.store(k_ptr + second_half_k_offsets, new_k_tile_2, mask=second_k_mask)
    else:
        new_q_tile_1 = q_tile_1 * cos_row + q_tile_2 * sin_row
        tl.store(q_ptr + first_half_q_offsets, new_q_tile_1, mask=first_q_mask)
        new_q_tile_2 = q_tile_2 * cos_row - q_tile_1 * sin_row
        tl.store(q_ptr + second_half_q_offsets, new_q_tile_2, mask=second_q_mask)

        new_k_tile_1 = k_tile_1 * cos_row + k_tile_2 * sin_row
        tl.store(k_ptr + first_half_k_offsets, new_k_tile_1, mask=first_k_mask)
        new_k_tile_2 = k_tile_2 * cos_row - k_tile_1 * sin_row
        tl.store(k_ptr + second_half_k_offsets, new_k_tile_2, mask=second_k_mask)


def rope_forward(q, k, cos, sin):
    q = q.transpose(1, 2)
    k = k.transpose(1, 2)

    batch_size, seq_len, n_q_head, head_dim = q.shape
    n_kv_head = k.shape[2]
    pad_hd = triton.next_power_of_2(head_dim)
    pad_n_q_head = triton.next_power_of_2(n_q_head)
    pad_n_kv_head = triton.next_power_of_2(n_kv_head)
    BLOCK_SIZE = max(pad_n_q_head, pad_n_kv_head)

    n_row = batch_size * seq_len

    q = q.contiguous()
    k = k.contiguous()
    cos = cos.contiguous()
    sin = sin.contiguous()

    _triton_rope[(n_row,)](
        q,
        q.stride(1),
        k,
        k.stride(1),
        cos,
        cos.stride(-2),
        sin,
        sin.stride(-2),
        seq_len,
        batch_size,
        n_q_head,
        n_kv_head,
        head_dim,
        pad_n_q_head,
        pad_n_kv_head,
        pad_hd,
        BLOCK_SIZE=BLOCK_SIZE,
        BACKWARD_PASS=False,
    )
    return q.transpose(1, 2), k.transpose(1, 2), cos, sin



##################################################################################################################################################


import torch

def test_rope_forward():
    # Define the test parameters
    batch_size = 16
    seq_len = 128
    n_q_head = 64
    n_kv_head = 64
    head_dim = 512

    # Create random input tensors
    q = torch.randn(batch_size, n_q_head, seq_len, head_dim, dtype=torch.float32, device='cuda')
    k = torch.randn(batch_size, n_kv_head, seq_len, head_dim, dtype=torch.float32, device='cuda')
    cos = torch.randn(seq_len, head_dim // 2, dtype=torch.float32, device='cuda')
    sin = torch.randn(seq_len, head_dim // 2, dtype=torch.float32, device='cuda')

    # Dictionary to store results for each test case
    results = {}

    # Test case 1: Forward pass
    q_out_1, k_out_1, cos_out_1, sin_out_1 = rope_forward(q, k, cos, sin)
    results['test_case_1'] = (q_out_1, k_out_1, cos_out_1, sin_out_1)

    # Test case 2: Backward pass
    q_out_2, k_out_2, cos_out_2, sin_out_2 = rope_forward(q, k, cos, sin)
    results['test_case_2'] = (q_out_2, k_out_2, cos_out_2, sin_out_2)

    return results

result_gold = test_rope_forward()

\end{lstlisting}
\end{mdframed}

\textbf{Our Converted PyTorch Reference Implementation:}
\begin{mdframed}[linewidth=1pt, linecolor=black, backgroundcolor=gray!5]
\begin{lstlisting}
import torch
import torch.nn as nn

class Model(nn.Module):
    """
    RoPE (Rotary Position Embedding) - PyTorch Reference Implementation

    Rotary Position Embedding applies a rotation to the query and key vectors
    based on their position in the sequence. This allows the model to naturally
    encode relative positions.

    Formula:
        For each position, split the embedding into two halves [x1, x2]
        Apply rotation: [x1*cos - x2*sin, x2*cos + x1*sin]

    Used in: LLaMA, GPT-J, GPT-NeoX, PaLM, and many modern LLMs
    """
    def __init__(self):
        super(Model, self).__init__()

    def forward(self, q, k, cos, sin):
        """
        Apply RoPE to query and key tensors.

        Args:
            q (torch.Tensor): Query tensor of shape (batch, n_heads, seq_len, head_dim)
            k (torch.Tensor): Key tensor of shape (batch, n_heads, seq_len, head_dim)
            cos (torch.Tensor): Cosine values of shape (seq_len, head_dim//2)
            sin (torch.Tensor): Sine values of shape (seq_len, head_dim//2)

        Returns:
            Tuple[torch.Tensor, torch.Tensor, torch.Tensor, torch.Tensor]:
                - q_rotated: Rotated query (batch, n_heads, seq_len, head_dim)
                - k_rotated: Rotated key (batch, n_heads, seq_len, head_dim)
                - cos: Cosine values (unchanged)
                - sin: Sine values (unchanged)
        """
        # Transpose to (batch, seq_len, n_heads, head_dim) for easier position-wise operation
        q = q.transpose(1, 2)
        k = k.transpose(1, 2)

        batch_size, seq_len, n_heads, head_dim = q.shape
        half_dim = head_dim // 2

        # Split into two halves along head_dim
        q1 = q[..., :half_dim]  # First half
        q2 = q[..., half_dim:]  # Second half
        k1 = k[..., :half_dim]
        k2 = k[..., half_dim:]

        # Reshape cos/sin for broadcasting: (seq_len, head_dim//2) -> (1, seq_len, 1, head_dim//2)
        cos = cos[:seq_len, :].unsqueeze(0).unsqueeze(2)
        sin = sin[:seq_len, :].unsqueeze(0).unsqueeze(2)

        # Apply rotation transformation
        # RoPE formula: rotate_half([x1, x2]) = [x1*cos - x2*sin, x2*cos + x1*sin]
        q_rotated = torch.cat([
            q1 * cos - q2 * sin,  # New first half
            q2 * cos + q1 * sin   # New second half
        ], dim=-1)

        k_rotated = torch.cat([
            k1 * cos - k2 * sin,
            k2 * cos + k1 * sin
        ], dim=-1)

        # Transpose back to (batch, n_heads, seq_len, head_dim)
        q_rotated = q_rotated.transpose(1, 2)
        k_rotated = k_rotated.transpose(1, 2)

        # Return cos/sin as well to match Triton interface
        return q_rotated, k_rotated, cos.squeeze(0).squeeze(1), sin.squeeze(0).squeeze(1)


BATCH_SIZE = 2
N_HEADS = 8
SEQ_LEN = 4
HEAD_DIM = 16

def get_inputs():
    """
    Generate test inputs for RoPE.

    Returns:
        List containing [q, k, cos, sin]:
            - q: Query tensor (batch, n_heads, seq_len, head_dim)
            - k: Key tensor (batch, n_heads, seq_len, head_dim)
            - cos: Cosine values (seq_len, head_dim//2)
            - sin: Sine values (seq_len, head_dim//2)
    """
    q = torch.randn(BATCH_SIZE, N_HEADS, SEQ_LEN, HEAD_DIM, dtype=torch.float32)
    k = torch.randn(BATCH_SIZE, N_HEADS, SEQ_LEN, HEAD_DIM, dtype=torch.float32)
    cos = torch.randn(SEQ_LEN, HEAD_DIM // 2, dtype=torch.float32)
    sin = torch.randn(SEQ_LEN, HEAD_DIM // 2, dtype=torch.float32)
    return [q, k, cos, sin]

def get_init_inputs():
    """
    Get initialization parameters for Model.

    Returns:
        Empty list (no initialization parameters needed)
    """
    return []
\end{lstlisting}
\end{mdframed}
